\documentclass{beamer}
\usetheme{Madrid}

\title[College Bus Surveillance]{Smart bus tracking and monitoring system using camera}
\subtitle{Zeroth Review Presentation}
\author[Group 6]{
\textbf{Group No: 6}\\
Akash V G\\
Fidha Fathima A\\
Mohammed Ramshad T S\\
Nithya Manoj
}
\date{}

\begin{document}

% Slide 1
\begin{frame}
    \titlepage
    \textbf{Guide name:Rakhi R}
\end{frame}

% Slide 2: Introduction
\begin{frame}{Introduction}
    \begin{itemize}
        \item {This project presents a mobile-based bus attendance system that uses face recognition and GPS to track students boarding and exiting the college bus.}
        \item {The solution is cost-effective, easy to deploy, and requires only a single smartphone without any additional hardware.}

    \end{itemize}
\end{frame}

% Slide 3: Objective
\begin{frame}{Objective}
    \begin{itemize}
        \item Develop an intelligent college bus surveillance and attendance system.
        \item Ensure accurate real-time student identification using face recognition.
        \item Track live bus location using GPS.
        \item Centralized dashboard for unified data storage and monitoring.
        \item Minimize hardware cost using smartphones instead of special devices.
    \end{itemize}
\end{frame}

% Slide 4: Existing System
\begin{frame}{Existing System}
    \begin{itemize}
        \item Manual attendance marking in buses leads to delays and errors.
        \item RFID or card-tap systems depend on students carrying cards.
        \item Additional hardware deployment increases cost and maintenance.
        \item No integrated real-time tracking and attendance solution available.
    \end{itemize}
\end{frame}

% Slide 5: Proposed System
\begin{frame}{Proposed System}
    \begin{itemize}
        \item Automatic attendance using facial recognition.
        \item Live tracking of college buses.
        \item Centralized storage of attendance and location data.
        \item Easy monitoring for administrators.
        \item Reduces manual work and errors.
    \end{itemize}
\end{frame}

% Slide 6: Base Paper Details
\begin{frame}{Base Paper Details}
    \textbf{Base Paper 1:}\\
    Bus Attendance System Through Facial Recognition by Implementing AI Integration\\
    Authors: Johnson C, Karthick M, Sakthivel P, Sasikala K, Santhi A (2023)

    \vspace{8pt}
    \textbf{Base Paper 2:}\\
    Smart College Bus Tracker: Real-Time Location and ETA Predictor\\
    Authors: M. Sowndharya, Yuvaraj K, Rahul S, Srichandiran S, Sriram S (2025)
\end{frame}

% Slide 7: Time Plan
\begin{frame}{Time Plan}
\begin{tabular}{|c|p{9cm}|}
\hline
\textbf{Week} & \textbf{Work} \\
\hline
1 & Requirement analysis, architecture, module identification, UI wireframing \\\hline
2 & DFD/ER, Flask backend setup, authentication, REST API \\\hline
3 & Web dashboard: admin login, student and bus management \\\hline
4 & Flutter app UI, camera preview, ML Kit detection \\\hline
5 & TensorFlow Lite recognition tuning and accuracy checks \\\hline
6 & Attendance workflow: GPS + timestamp, offline sync \\\hline
7 & Real-time GPS tracking, map integration (Google Maps/OSM) \\\hline
8 & Integration, testing, low-light evaluation, optimization \\\hline
9 & UI polishing, security, performance tuning \\\hline
10 & LaTeX report preparation \\
\hline
\end{tabular}
\end{frame}

% Slide 8: References
\begin{frame}{References}
    \begin{enumerate}
1. IJRPR, Vol. 8, Issue 11, 2023 – Facial Recognition Attendance System.\\
2. IJSRED, Vol. 8, Issue 3, 2025 – Smart Bus Tracking and ETA Predictor.\\
3. Google ML Kit Documentation.\\
4. TensorFlow Lite Developer Guide.\\
5. Google Maps & OpenStreetMap API Documentation.
    \end{enumerate}
\end{frame}

\end{document}
\textbf{\textit{\textbf{}}}